\section{Attachment to the full development and the observing geodesic}
\label{attach:section}

Fix the background datum $d$ and the branch seed $w$ of
Corollary~\ref{test:branch-seed}. Use the separately prepared branch
coordinates in which $v\to\infty$ and $u\to u_*<\infty$ along
$\gamma_w$. Write $p_T=\gamma_w\cap\{v=T\}$, and denote the two
complete corrected data of Proposition~\ref{prep:signed-entry} by
$d_T^\pm$. The exact interior metrics of
Proposition~\ref{energy:signed-experiment} are denoted by $g_T^\pm$.
This section places their finite observations in the corresponding
full maximal globally hyperbolic developments and supplies observing
seeds in the original open seed set.

After the fixed derivative orders and the small losses have been
chosen as in Proposition~\ref{energy:signed-experiment}, there are
$c>0$ and constants independent of $T$ such that
\begin{equation}\label{attach:relative-smallness}
 \lVert H_T^\pm\rVert_{C^{30}_{\mathrm{norm}}}
 \le \varepsilon_T,
 \qquad \varepsilon_T=C\exp\bigl(- (\kappa/2+c)T\bigr).
\end{equation}
Here the norm means ordinary derivatives of the normalized metric
arrays. Indeed, the packet has exponent $-7\kappa/8$ before its
allocated frequency loss, while the largest individual error has
exponent $-3\kappa/4$ before the allocated rate and finite-jet losses.
These losses can be chosen strictly smaller than the gap to
$-\kappa/2$. A further fixed-order factor $e^{o(T)}$ can always be
absorbed after decreasing $c$. In particular
$T\varepsilon_T\to0$ and
$\varepsilon_T/\sqrt{a(p_T)}\to0$.

\subsection{A Cauchy development of the entire corrected datum}\label{attach:whole-bridge-section}

We join three vacuum regions: a lower development containing the
complete bridge and a past collar, the selected late exterior
evolution, and the finite deep evolution. The lower region must retain
every initial spatial radius, including the corrected exterior tail.
Every discarded face above it must be a future exit, so that a past
causal curve from the experiment can reach the lower region. These
are the two geometric requirements behind the attachment.

The finite bridge lemma supplies the lower region, the delayed exterior
partial slices, and enlarged fixed matching bands
\cite[\CompanionRef{foundation}{cmp:finite-range}]{Foundation}.
The comparison-development proof supplies their fixed-time wave-map
identifications and product collars
\cite[\CompanionRef{foundation}{cmp:development}]{Foundation}; the
packet-attachment proof supplies the strict future cuts and the order
of the joins
\cite[\CompanionRef{foundation}{pk2:attachment-proposition}]{Foundation}.
We first identify the actual solutions on these collars. We then choose
the retained regions and prove that the complete bridge is Cauchy before
invoking maximality. The verification below uses the present signed
estimates, whose amplitude differs from the one in that source.

\begin{proposition}[Whole-bridge attachment]\label{attach:development}
For all sufficiently large finite $T$, there is a smooth time-oriented
vacuum development $\mathcal N_T^\pm$ of the entire datum $d_T^\pm$
with $\Sigma$ as a Cauchy hypersurface. It contains the constructed
finite observation, restricted within its construction margins, and
an optical-coordinate box about $p_T$, of fixed sufficiently small
coordinate radii independent of $T$, in
$I^+_{\mathcal N_T^\pm}(\Sigma)$.
It admits a time-orientation-preserving open isometric embedding into
$\mathcal M_{d_T^\pm}$ that preserves the initial-data embedding.
The full physical constraint-correction tail is part of its initial
data.
\end{proposition}

\begin{proof}
Fix one sign. Perform each identification on enlarged solution regions;
the retained regions will be chosen after these identifications.

\emph{The full lower block.}
The data converge to $d$ in $p_{10}$, are smooth, and solve the
vacuum constraints. Hence the finite bridge lemma applies with the
same fixed depths, time bands, and past collar. Its future artificial
faces are strictly future-exiting for causal curves. At bounded
radius its bounds are finite-time Cauchy bounds on compact causal
shadows; on the ends they are uniformly local symbol bounds, with
bounded causal coordinate speeds. Its smallness hypothesis is precisely
the fixed $p_{10}$ hypothesis, and no smallness at all smooth orders
is needed.

Let $t_T$ be the exterior launch time, let $C_e$ be the exterior
center, and let $s_* =\sup|C_e'|<1-2\beta$. The protected ball from
Proposition~\ref{prep:signed-entry} contains the fixed-origin ball
of radius
\begin{equation}\label{attach:protected-radius}
 (1-2\beta)t_T-|C_e(t_T)|
 \ge (1-2\beta-s_*)t_T-|C_e(0)|\longrightarrow\infty.
\end{equation}
Thus every fixed compact causal shadow used in the lower joins has
unchanged data for large $T$. Domain of dependence identifies the
bounded central pieces, the bent compact part of the selected
partial slice, and the required pieces on the unselected end with
the background ones. On the horizontal far part of the selected
partial slice $t_e=0$, however, the lower development carries the
actual datum $d_T^\pm$, including its entire tail. This is also the
initial datum of the selected exterior evolution. Equality with the
old datum is asserted only in the protected regions.

\emph{The fixed-time exterior identification.}
Choose the fixed matching time $t_a>0$ and an enlarged band around
it supplied by the comparison-development construction. Both the
lower block and the selected exterior solution evolve the same
geometric data on the partial slice. Solve the identifying wave map $\Psi$
from the exterior solution to the lower solution with the slice
identification as its value and the matching future unit normal as
its normal derivative. The background identification may be used as
the reference coordinate map. On the fixed slab from $t_e=0$ through
the matching band, the map remains a diffeomorphism on the smaller
region.

Here is the quantitative input for this last assertion. The source
construction gives metric $H^{10}$ and second-form $H^9$ control on
compact shadows and the corresponding uniformly local end control.
The differentiated wave-map equation at map order ten uses
Christoffel derivatives through nine. In a bootstrap with bounded
low map derivatives and inverse Jacobian, tame product and composition
estimates bound its current-order coefficient by the uniformly bounded
low metric and map norms. The known higher coefficient derivatives
through the stated finite orders multiply lower map derivatives; their
norms are controlled by the same fixed-time package. Thus the energy
of the difference from the background map obeys an inequality of the
form
\begin{equation}\label{attach:map-estimate}
 E_{10}(s)\le C\eta_T+C\int_0^s E_{10}(\sigma)\,d\sigma,
 \qquad 0\le s\le t_a+\delta_a,
 \qquad \eta_T\longrightarrow0,
\end{equation}
where $C$ and $\delta_a>0$ do not depend on the growing outer cutoff.
On the ends the estimate is taken as a supremum of energies in
uniformly local charts, so its coefficient does not grow with their
number or total volume. Uniqueness identifies the maps on chart
overlaps. Fixed-time Gronwall, Sobolev embedding, and the strict
spatial and deeper margins give the required $C^1$ diffeomorphism
control. To see global injectivity on the used region, compose with
the inverse background identification. The resulting map $\widehat\Psi$ is
uniformly close to the identity in both $C^0$ and $C^1$. Two points
with the same image must therefore be close in the reference
bounded-geometry metric. They lie in one enlarged convex chart,
where integration of $D\widehat\Psi=\Id+o(1)$ along their joining segment
rules out a collision. The inverse function theorem supplies the
smooth inverse on the image. We also obtain at least $H^9$ control
after straightening the joining
collar. This is exactly the finite-order map estimate in the cited
comparison-development proof. Late high-order bounds of size
$e^{o(T)}$ are not used as its top Gronwall coefficient.

Write $g_E,g_L$ for the exterior and lower metrics.
Set $h=\Psi^*g_L$ and now hold this
smooth metric fixed as the target. The wave-map identity is
$\operatorname{tr}_{g_E}(\Gamma(g_E)-\Gamma(h))=0$.
Thus $g_E$ and $h$ are vacuum metrics in the same target-$h$
reduced system. Slice and future-normal matching give equal initial
metric values and tangential derivatives. Equality of the second
fundamental forms fixes their normal spatial-metric derivatives.
More explicitly, let $B_{\mu\nu}=\partial_0(g_E-h)_{\mu\nu}$
on the slice, so $B_{ij}=0$. Lowering the relative gauge vector
with the common initial metric gives
\[
 \mathcal C_i=h^{00}B_{0i},\qquad
 \mathcal C_0=\tfrac12h^{00}B_{00}.
\]
Since $\mathcal C=0$ and $h^{00}\ne0$, the remaining normal
derivatives agree as well.
The two metrics have identical full first reduced jets. Smooth
local uniqueness, with the strict outflow boundaries of the
enlarged domain, proves that the map is an isometry throughout the
used matching band. This uniqueness argument is for each fixed $T$;
its higher smooth norms do not enter the preceding uniform map
estimate. The small coordinate straightening can be cut
off farther back inside the available margin. The compact joins use
the unchanged background identification.

\emph{Entry into the deep domain.}
The selected exterior exists through $t_T+C_2$, with a shrinking
outer cutoff and strict deeper exit boundaries, by
Proposition~\ref{prep:signed-entry}. The used interval of
$\mathcal S=\{t=0\}$ and its collar lie inside these margins: its
selected late endpoint has bounded co-radius and exterior time
$t_e=t_T+O(1)$, whereas its middle and unselected parts are supplied
by the unchanged lower pieces. The spacelike conormal of this entry
surface retains its strict sign under the low-order comparison.
The deep evolution uses exactly the induced spacetime jets furnished
by this exterior and the same target-$g$ wave-map gauge. Consequently
local reduced uniqueness identifies their actual entry collars.
An approximate packet match alone would not suffice for this step.

The three metrics now agree under actual collar isometries. It remains
to choose which portions to keep and to show that a past causal curve
from any retained future point reaches the complete lower block.

\emph{Retained domains and their endpoints.}
Let $r=|y|$ be the selected-end co-radius. Beyond the fixed join only
the used exterior region is retained, with outer radius
$R_{\mathrm{out}}(t_e)$ and terminal time $t_{\mathrm{end}}$.
The speed chosen in Proposition~\ref{prep:signed-entry} satisfies
\begin{equation}\label{attach:outer-sign}
 \frac{d}{dt_e}\bigl(r-R_{\mathrm{out}}(t_e)\bigr)
 \ge V_{\mathrm{out}}-V_{\max}>0
\end{equation}
on future causal curves near the active outer cut. The retained
selected-end set near the matching face is
\begin{equation}\label{attach:union}
 \{t_e<t_{\mathrm{end}}\}\cap
 \bigl(\{t_e<t_a\}\ \cup\
       \{r<R_{\mathrm{out}}(t_e)\}\bigr),
\end{equation}
with the lower entry-surface cut imposed where active. The first
member of the union contains all radii. In particular no physical
initial tail is replaced by background data beyond the computing
cutoff. Figure~\ref{attach:union-figure} depicts this local condition.

\begin{figure}[ht]
\centering
\begin{tikzpicture}[x=1.15cm,y=0.9cm,>=Stealth,font=\small]
 \fill[blue!9] (0,0) rectangle (5.5,1.05);
 \fill[blue!9] (0,1.05)--(4.15,1.05)--(3.1,3.15)--(0,3.15)--cycle;
 \fill[orange!18] (0.15,0.82) rectangle (3.8,1.28);
 \draw[->] (0,-0.12)--(5.9,-0.12) node[right] {$r$};
 \draw[->] (-0.12,0)--(-0.12,3.55) node[above] {$t_e$};
 \draw[densely dashed] (0,1.05)--(5.5,1.05);
 \node[left] at (-0.12,1.05) {$t_a$};
 \draw[thick,blue!65!black] (0,1.05)--(4.15,1.05);
 \draw[thick,dashed] (4.15,1.05)--(3.1,3.15);
 \node[align=center] at (1.65,2.4) {retained future end\\$r<R_{\mathrm{out}}(t_e)$};
 \node[align=center] at (4.8,2.6) {omitted};
 \node[anchor=west,rotate=-63] at (3.7,1.95) {$R_{\mathrm{out}}$};
 \node at (2.5,0.4) {full lower end: every radius};
 \node[anchor=west] at (0.22,1.57) {matching collar};
 \draw[->] (0.95,1.51)--(0.95,1.24);
 \draw[->,thick] (2.9,1.8)--(2.7,0.6);
 \node[right] at (2.9,1.45) {past};
 \draw[->,blue!65!black] (5.15,0.4)--(5.75,0.4);
\end{tikzpicture}
\caption{The local exterior condition in Equation~\eqref{attach:union}.
The lower part extends to all radii; only the added future part has
an outer cut. The dashed outer edge is omitted. The matching segment
has a genuine collar. Terminal and entry-surface cuts, imposed
separately, are suppressed in this schematic.}
\label{attach:union-figure}
\end{figure}

For a past curve initially in the added future part,
$f=r-R_{\mathrm{out}}$ is nonincreasing as long as that cut is
active. It therefore keeps the strictly negative gap it had at its
starting point until reaching $t_e=t_a$. It reaches an actual
matching collar there, then the unrestricted lower block. A curve
already in $t_e<t_a$ stays in that union member. In particular
$(t_e,f)=(t_a,0)$ cannot be a past endpoint from the added part.
The unmatched part of the $t_a$ face is a future exit and is omitted.

For clarity, the retained deep domain has strict future inequalities
\begin{equation}\label{attach:deep-domain}
 \begin{split}
 t&<L_T:=\frac{T+u_*}{2}+c_0,\qquad t>0,\\
 F_0:=u+\frac{c_0t}{T}&<u_2,
 \qquad F_1:=v+\frac{c_0t}{T}<T+c_2,
 \end{split}
\end{equation}
where $u_2>u_*$ and $c_2>0$ are fixed, and $c_0>0$ was chosen
small relative to these margins. The lower face $t=0$ is attached.
Relative cone control from Equation~\eqref{attach:relative-smallness}
keeps $t$ temporal and gives, for future causal vectors normalized
by $dt=1$,
\begin{equation}\label{attach:deep-sign}
 du,dv\ge-C\varepsilon_T,
 \qquad dF_i\ge \frac{c_0}{T}-C\varepsilon_T>0.
\end{equation}
Indeed, after division of the causal inequality by $a$, the null
components are $du,dv$ and the angular velocity is
$(d\theta-b\,du)/\sqrt a$. On the section $du+dv=2$, the
background cone is a bounded set and has $du,dv\ge0$;
a small relative perturbation enlarges these lower bounds by at
most $C\varepsilon_T$. This also bounds the angular velocity on
that section.

If a past curve starts at $q$ inside this domain, then along it
$F_i\le F_i(q)$. Thus its gaps from both omitted sides stay
strictly positive, rather than merely nonnegative. At a limiting
point on $t=0$ these same gaps put it strictly inside the used entry
interval. This interval and its collar were chosen strictly inside
the exterior interval available through $t_T+C_2$, including its
outer and deeper margins. Such a limit is therefore in the actual
identified collar. It cannot be an entry/side or entry/outer-cut
corner. The terminal time inequality has the same past-preservation
property. The other active temporal and inner cuts in the lower
matches have the source construction's strict future-exit signs.

At an intersection of cuts, the retained prescription is a finite
union or intersection of these past-preserved strict sublevels.
For a union, one retains the member occupied by the curve until it
enters an unrestricted lower member. For an intersection, every
active gap is preserved. This locates any finite past endpoint in
an ordinary retained chart, an actual matching collar, or the full
lower block. No continuation across an omitted face is used.
For each fixed $T$ the added coordinate ranges are finite and the
metric is smooth and nondegenerate on the enlarged ranges. A causal
curve parametrized by the relevant temporal coordinate has bounded
coordinate speeds there, hence a limit at a finite endpoint. The
gap argument just given places this limit in a retained chart,
where causal continuation is possible. In the lower end pieces,
the fixed-time speed bound rules out escape to spatial infinity.

\emph{The gluing topology and the Cauchy property.}
Straighten each matching collar as a product with the matching
hypersurface. Retain the lower side from the lower solution and the
upper side from the added solution, with only one copy of the
matching face. The interiors of these pieces are disjoint. Remove
unused future faces and their edges before taking the quotient;
the matching faces are then closed relative to their respective
retained pieces. The identifications are homeomorphisms on those
closed faces and isometries on their collars. A finite gluing of
this kind is Hausdorff: saturating a closed subset adds its closed
image in each matching face, so the quotient map is closed; its
fibres are finite, and distinct fibres can be separated in the
Hausdorff pieces and then by saturated neighborhoods. The product
collars give the smooth charts across the faces. At intersecting
joining prescriptions first join the two lower charts, then cross
$\mathcal S$; the uniqueness identifications agree on the overlap.
Thus there is no duplicate open continuation of a single face.
Second countability follows from the finitely many countable atlases.

Every past-inextendible causal curve starting above $\Sigma$ now
reaches the lower block and the complete bridge: it cannot stop at
an omitted face, at a finite interior endpoint, or at infinity in
bounded time. The joining faces are spacelike and have a fixed
crossing orientation, so the finitely many joining layers cannot
be revisited indefinitely. The lower block has the Cauchy control
of the finite bridge lemma. A curve in the short past collar reaches
$\Sigma$ to the future by that same lemma. No causal curve crosses
$\Sigma$ twice, by lower-block Cauchyness and the one-way signs of
the subsequent joining faces. Hence every inextendible causal curve
meets $\Sigma$ exactly once. The glued metric is smooth vacuum and
induces exactly $(h_T^\pm,K_T^\pm)$ on this full hypersurface.
It is therefore a globally hyperbolic development in the sense of
\cite[Definition~2.1]{Sbierski2016}. The maximal-development theorem
\cite[Definition~2.2 and Theorem~2.8]{Sbierski2016} gives the claimed
initial-data-preserving isometric embedding into
$\mathcal M_{d_T^\pm}$; an embedding of equal-dimensional manifolds
is open. At $p_T$, the gap from the time top is at least $c_0$, and the
two side gaps are at least $u_2-u_*-c_0/2-O(T^{-1})$ and
$c_2-c_0/2-O(T^{-1})$. These are fixed positive numbers.
The derivatives of the cut functions in $(u,v)$ are uniformly
bounded. A fixed sufficiently small coordinate box therefore lies
inside the retained deep domain for every large $T$. It has positive
time above the spacelike entry, so it lies in $I^+(\Sigma)$.
\end{proof}


\begin{remark}[Geometric hypotheses of the attachment]
\label{attach:geometric-contract}
The proof of Proposition~\ref{attach:development} uses the smallness
of the interior change only to keep the time levels spacelike and
the two tilted side functions strictly increasing on future causal
vectors. Thus the same attachment proof applies to a finite vacuum
experiment with exactly matched entry Cauchy jets, complete corrected
data and the stated protected region, provided it has these temporal
and strict future-outflow inequalities and the fixed matching margins.
The finite-time exterior identification still requires the low-order
closeness used in Equation~\eqref{attach:map-estimate}. The stronger
estimate $H=o(\sqrt a)$ is first used in the adaptive observer
construction below; it is not a geometric gluing hypothesis.
\end{remark}

\subsection{The absolute canonical momentum on the unit shell}

We will prescribe the perturbed observer at its terminal point and solve
backwards. First we construct its geodesic vector field with $v$ as
time and estimate its difference from the background field.
Lemma~\ref{bg:branch} applies without any
assumption on parallel-frame curvature. With $U=\dot\gamma_w$, its conclusions are
\begin{equation}\label{attach:branch-variables}
 \Xi=\dot\theta-b\dot u,\qquad p=a\dot v\asymp1,
 \qquad P_A=\gamma_{AB}\Xi^B,
 \qquad |P|+|\Xi|+|\dot u|\le C,
 \qquad a(p_T)\asymp e^{-\kappa T}.
\end{equation}
The angular path has a limit. In canonical cotangent coordinates
$\pi=\pi_u\,du+z\,dv+P_A\,d\theta^A$, where $z=\pi_v$, the sign
of the shifted momentum is
\begin{equation}\label{attach:shifted-momentum}
 p=-(\pi_u+b^AP_A).
\end{equation}
The background unit shell and its $v$-time Hamiltonian are
\begin{equation}\label{attach:background-shell}
 2pz+a(1+\gamma^{AB}P_AP_B)=0,
 \qquad \mathcal H_0=-z_0
   =\frac{a(1+\gamma^{AB}P_AP_B)}{2p}.
\end{equation}
The phase variables for this Hamiltonian are
$Y=(u,\theta,\pi_u,P)$; $v$ is the time variable.

Along the background track $z_0=O(a)$, whereas its perturbation may be larger
than $a$. We therefore estimate $z=\pi_v$ itself, rather than $z/a$.
The next two lemmas put the background track in a fixed state tube and
construct a future-pointing perturbed root there. Their output is a
$v$-time vector field whose absolute difference from the background is
at most $\varepsilon_T e^{o(T)}$; this is the error that will be
integrated when returning the observer to the initial surface.

\begin{lemma}[A fixed canonical tube]\label{attach:canonical-tube}
There are a fixed late $v_0$ and a fixed bounded open canonical
state set $\mathcal K$ such that the late track stays a fixed
positive distance from $\partial\mathcal K$, its positions lie
in one angular coordinate chart, and $p(v,Y)\ge p_*>0$ for every
$v\ge v_0$ and $Y\in\mathcal K$. On a slightly enlarged such set,
the background canonical vector field $V_0$ has
\begin{equation}\label{attach:integrable-rate}
 |V_0(v,Y)|+|D_YV_0(v,Y)|
 \le C e^{-\kappa v}e^{o(v)}=:\Lambda(v),
 \qquad \int_{v_0}^{\infty}\Lambda(v)\,dv<\infty.
\end{equation}
The expression $e^{o(v)}$ here uses the current endpoint $v$, not
the terminal experiment parameter $T$.
\end{lemma}

\begin{proof}
The state momenta and $u$ are bounded by
Equation~\eqref{attach:branch-variables} and
Equation~\eqref{attach:shifted-momentum}. Choose $v_0$ so that the
angular track is in an inner subchart about its limit, with a fixed
coordinate margin. The bounded first derivatives of $b$ in
Proposition~\ref{bg:geometry} imply, on a fixed bounded momentum
range and at the same $v$,
\[
 |p(Y)-p(Y_0(v))|
 \le C\bigl(|u-u_0(v)|+|\theta-\theta_0(v)|
              +|\pi_u-\pi_{u,0}(v)|+|P-P_0(v)|\bigr),
\]
with $C$ independent of $v$. Choose $r_0>0$ smaller than the angular
margin and than $\inf p(Y_0(v))/(2C)$. On every radius-$r_0$
neighborhood of $Y_0(v)$, the momentum $p$ then stays uniformly positive.
This is a pointwise tube choice and entails no multiplication of
$C$ by the length of the late interval.

Twice differentiating the Hamiltonian in
Equation~\eqref{attach:background-shell} uses a fixed finite number
of ordinary background jets and bounded inverse powers of $p$.
Every derivative of $a$ retains $a$, and all remaining fixed jets
cost $e^{o(v)}$ on a bounded $u$ interval by
Proposition~\ref{bg:geometry}. Thus both the Hamiltonian vector field
and its first state derivative retain $a\asymp e^{-\kappa v}$,
which proves Equation~\eqref{attach:integrable-rate} on these
constant-radius neighborhoods. In particular $Y_0(v)$ has a limit
$Y_\infty$. Increase $v_0$ until $|Y_0(v)-Y_\infty|<r_0/8$
throughout the tail. A fixed ball $\mathcal K$ of radius $r_0/4$
about $Y_\infty$ then contains the track with margin $r_0/8$,
and even its radius-$r_0/2$ enlargement lies in the already
controlled radius-$r_0$ neighborhood of $Y_0(v)$ for every $v$.
This gives a state domain independent of both $v$ and $T$.
\end{proof}

Near the finite experiment set
\begin{equation}\label{attach:isometry}
 S_T^\pm=g^{-1}g_T^\pm,
 \qquad I_T^\pm=(S_T^\pm)^{-1/2}.
\end{equation}
The inverse square root is the real analytic branch at the identity,
defined by its convergent power series. Since $S_T^\pm$ is
$g$-self-adjoint, so is $I_T^\pm$, and the two commute. Consequently
\[
 g_T^\pm(I_T^\pm V,I_T^\pm W)
 =g(S_T^\pm I_T^\pm V,I_T^\pm W)=g(V,W).
\]
It is future-preserving by continuity from the identity. In a
background normalized orthonormal angular frame, the background
Gram matrix is constant. Differentiating the power series and using
Equation~\eqref{attach:relative-smallness} therefore gives, at the
fixed orders used below,
\begin{equation}\label{attach:isometry-bound}
 \lVert I_T^\pm-\Id\rVert_{C^{30}_{\mathrm{norm}}}
 +\lVert (g_T^\pm)^{-1}-g^{-1}\rVert_{C^{30}_{\mathrm{norm}}}
 \le \varepsilon_T e^{o(T)}.
\end{equation}
The inverse-metric norm refers to coefficients in the normalized
vector basis, rather than coordinate inverse components. We decrease
$c$ if needed to absorb these fixed-order factors when using
Equation~\eqref{attach:isometry-bound} as a terminal bound.

\begin{lemma}[The absolute future root]\label{attach:absolute-shell}
On the fixed canonical tube of Lemma~\ref{attach:canonical-tube},
intersected with the actual deep domain, the new unit shell has a
simple root $z=z_T^\pm(v,Y)$ close to $z_0(v,Y)$, with
\begin{equation}\label{attach:root-bound}
 \lVert z_T^\pm-z_0\rVert_{C^2_Y}
 \le \varepsilon_T^{\mathrm{sh}},\qquad
 \varepsilon_T^{\mathrm{sh}}=\varepsilon_T e^{o(T)}.
\end{equation}
The associated canonical vector field satisfies
$|V_T^\pm-V_0|\le C\varepsilon_T^{\mathrm{sh}}$. On this root $dv/d\tau>0$.
The estimates concern $\pi_v$ itself.
\end{lemma}

\begin{proof}
Put $g'=g_T^\pm$. Choose a background angular orthonormal frame
$e_A$ in the fixed chart and write $P_{(A)}=\pi(e_A)$. Let
$Q^{\alpha\beta}$ be the symmetric coefficients of
$g'^{-1}-g^{-1}$ in the vector basis
$(a^{-1/2}L_0,a^{-1/2}L_1,e_A)$. They and the required ordinary jets
are $O(\varepsilon_T e^{o(T)})$. Multiplication of the new shell by
$a$ gives the exact equation
\begin{equation}\label{attach:multiplied-shell}
 \mathcal F_T(v,Y,z)
 =2pz+a(1+\gamma^{AB}P_AP_B)+\Delta_T(v,Y,z)=0,
 \qquad \Delta_T=a(g'^{-1}-g^{-1})(\pi,\pi).
\end{equation}
Keeping all powers of $a$, its last term is the polynomial
\begin{align}\label{attach:shell-polynomial}
 \Delta_T={}&Q^{00}p^2-2Q^{01}pz+Q^{11}z^2
       -2\sqrt a\,Q^{0A}pP_{(A)}
       +2\sqrt a\,Q^{1A}zP_{(A)}\notag\\
       &\hspace{32mm}+aQ^{AB}P_{(A)}P_{(B)}.
\end{align}
This follows from $\pi(L_0)=-p$ and $\pi(L_1)=z$.
It displays the coefficients of sizes $\varepsilon_T$,
$\sqrt a\,\varepsilon_T$, and $a\varepsilon_T$ before their
fixed-order losses. In particular the coefficient of $z^2$ is
allowed to be nonzero.

On bounded $(Y,z)$ sets, ordinary and state differentiation of this
polynomial gives
\begin{equation}\label{attach:polynomial-jets}
 |\partial_v^{j_0}D_Y^{j_1}\partial_z^{j_2}\Delta_T|
 \le C_j\varepsilon_T e^{o(T)}
 \quad (j_0+j_1+j_2\le3).
\end{equation}
A derivative of $\sqrt a$ or $a$ retains that factor; a derivative
of $p$ uses $b$ and its ordinary derivatives; the angular
orthonormalization uses only fixed-order background jets. Thus
frame differentiation introduces no additional inverse power of
$a$. The two null covector factors have already been cancelled by
the external $a$ in Equation~\eqref{attach:multiplied-shell}.

Fix a bounded $z$ interval containing $z_0$ with a fixed margin.
For large $T$, Equation~\eqref{attach:polynomial-jets} implies
\begin{equation}\label{attach:simple-root}
 \partial_z\mathcal F_T=2p+\partial_z\Delta_T\ge p_*>0.
\end{equation}
The values at $z_0\pm C\varepsilon_T^{\mathrm{sh}}$ have opposite signs for $C$ large,
so there is a unique root in this neighborhood. Implicit
differentiation, using the denominator in
Equation~\eqref{attach:simple-root}, proves
Equation~\eqref{attach:root-bound}; any finite background jet factors
in the resulting products are absorbed in the displayed $e^{o(T)}$.
The reduced Hamiltonian is $\mathcal H_T^\pm=-z_T^\pm$, which gives
the asserted vector-field bound. This reasoning does not claim that
$z_T^\pm-z_0$ is small relative to $a$.

To check both the parametrization and its time orientation, let
$\mathcal H_{\mathrm{aff}}=(g'^{-1}(\pi,\pi)+1)/2$ be the affine geodesic
Hamiltonian. Since $a$ is independent of $z$,
\begin{equation}\label{attach:future-root}
 \partial_z\mathcal F_T
   =2a\,\partial_z\mathcal H_{\mathrm{aff}}
   =2a\,\frac{dv}{d\tau}.
\end{equation}
Equation~\eqref{attach:simple-root} is therefore exactly the
positive-$v$ branch condition. For the time orientation, the
function $t=(u+v)/2$ remains temporal under the small relative
metric perturbation. At fixed position and angular momentum,
$\partial_{\pi_u}p=-1$, so the other momentum derivative gives
\[
 4a\frac{dt}{d\tau}
 =\partial_{\pi_u}\mathcal F_T+\partial_z\mathcal F_T
 =2p-2z+\partial_{\pi_u}\Delta_T+\partial_z\Delta_T
 \ge 2p_*-C\varepsilon_T^{\mathrm{sh}}>0.
\]
Here $z_0\le0$, $|z-z_0|\le\varepsilon_T^{\mathrm{sh}}$, and
Equation~\eqref{attach:polynomial-jets} controls both error
derivatives. Thus the root is future pointing using only absolute
root closeness. The other canonical equations,
divided by $dv/d\tau$, are those of $-z_T^\pm$ by implicit
differentiation of the shell. This remains true in the presence of
the quadratic term in $z$.
\end{proof}

\subsection{Returning the observer to the initial surface}

The shell construction controls the geodesic vector field on a fixed
state tube. We now choose the terminal unit tangent by the isometry
$I_T^\pm$, follow that state backwards within the actual attached
domain, and compare the resulting observer frame with the test norm.

\begin{proposition}[Adaptive observer and test-frame comparison]
\label{attach:observer}
At $p_T$ put $U_T^\pm=I_T^\pm U$. There are seeds
$w_T^\pm\to w$ on $\Sigma$ whose future unit geodesics in
$\mathcal M_{d_T^\pm}$ reach $p_T$ with tangent $U_T^\pm$, at finite
positive proper times strictly before their full-MGHD future
durations. In particular $w_T^\pm$ belong to the original open
seed set for all sufficiently large $T$.

Choose an angular orthonormal frame $e_A$ near $p_T$, including the
packet polarization direction at the center, and set
\begin{equation}\label{attach:adapted-frame}
 E_0^\pm=I_T^\pm L_0,\qquad
 E_1^\pm=a^{-1}I_T^\pm L_1,\qquad
 E_A^\pm=I_T^\pm e_A.
\end{equation}
At $p_T$ the Euclidean norm of components in this frame and the
parallel positive test norm $k_{d_T^\pm,w_T^\pm}$ are uniformly
equivalent, with constants independent of $T$. Their induced
operator norms are uniformly equivalent as well.
\end{proposition}

\begin{proof}
\emph{Terminal state.}
The normalized components of $U$ at $p_T$ are
\[
 \sqrt{a_p}\,\dot u,\qquad p/\sqrt{a_p},\qquad \Xi^{(A)},
 \qquad a_p=a(p_T),
\]
and hence have size $O(a_p^{-1/2})$. The new covector is
$\pi'=g_T^\pm(I_T^\pm U,\cdot)
       =g(U,(I_T^\pm)^{-1}\cdot)$.
Equation~\eqref{attach:isometry-bound} and the coordinate identities
$\partial_u=\sqrt a\,\hat e_0-b^A\partial_A$ and
$\partial_v=\sqrt a\,\hat e_1$ give
\begin{equation}\label{attach:terminal-state}
 |Y_T^\pm(T)-Y_0(T)|
 \le C\varepsilon_T/\sqrt{a_p}=o(1),
 \qquad |\pi'_v-\pi_v|\le C\varepsilon_T,
 \qquad dv(U_T^\pm)=\frac{p+O(\varepsilon_T)}{a_p}>0.
\end{equation}
The terminal positions agree exactly, and the angular coordinate
and orthonormal bases have bounded conversion matrices. The new covector
satisfies the exact unit shell. It lies in the fixed bounded root
interval and in the fixed state tube, so uniqueness in
Lemma~\ref{attach:absolute-shell} places it on that lemma's future
root. The loss $a_p^{-1/2}$ in the first estimate of
Equation~\eqref{attach:terminal-state} is paid once, at this terminal
conversion.

\emph{Backward comparison on the actual domain.}
Increase the fixed $v_0$ if necessary so that
$(v_0+\inf u)/2$ has a positive margin above zero. Choose the state
tube with a smaller concentric tube around the background track,
leaving a fixed boundary margin. For $v_0\le v\le T$ the background
spacetime track has the uniform gaps
\begin{align}\label{attach:track-margins}
 u_2-u(v)-\frac{c_0t(v)}{T}
   &\ge u_2-u_* -c_0/2-O(T^{-1}),\notag\\
 T+c_2-v-\frac{c_0t(v)}{T}
   &\ge c_2-c_0/2-O(T^{-1}),\\
 L_T-t(v)&\ge c_0,
 \qquad t(v)\ge t(v_0)>0.\notag
\end{align}
Here $u$ and $v$ increase along the background future timelike
track. All four right-hand margins are positive after the original
choice of $c_0$ and for large $T$.

Solve backward from the terminal state only as long as the solution
lies both in the fixed state tube and in the actual domain
of Equation~\eqref{attach:deep-domain}, stopping no later than
$v_0$. On any such provisional interval split the field difference
as
\[
 V_T^\pm(v,Y_T^\pm)-V_0(v,Y_0)
 =\bigl(V_T^\pm(v,Y_T^\pm)-V_0(v,Y_T^\pm)\bigr)
  +\bigl(V_0(v,Y_T^\pm)-V_0(v,Y_0)\bigr).
\]
The first term is bounded by $C\varepsilon_T^{\mathrm{sh}}$ and the second by
$\Lambda(v)|Y_T^\pm-Y_0|$. Only the background Lipschitz constant
is integrated. Backward Gronwall gives, uniformly in the provisional
stopping point,
\begin{equation}\label{attach:backward-error}
 \sup_{v\le s\le T}|Y_T^\pm(s)-Y_0(s)|
 \le \exp\!\left(\int_{v_0}^\infty\Lambda(s)\,ds\right)
       C\left(\frac{\varepsilon_T}{\sqrt{a_p}}+T\varepsilon_T^{\mathrm{sh}}\right)
 =:\delta_T\longrightarrow0.
\end{equation}
For large $T$, this is smaller than half every state-tube margin
and every spacetime margin in Equation~\eqref{attach:track-margins}.
It therefore excludes a first stop at the state boundary, either
tilted side, the top time cut, or $t=0$. At a finite remaining
interior endpoint the smooth finite-$T$ canonical ODE continues
inside its actual domain. Thus the solution reaches $v=v_0$.
The argument has not extended a perturbed vector field beyond the
constructed spacetime to justify its own estimate.

\emph{Return through the earlier compact segment.}
From the fixed section $v=v_0$ to $\Sigma$ the background observer
runs along a fixed compact segment. Choose a compact flow
neighborhood with strict margins in the entry and lower matching
collars; the crossings of $\mathcal S$ and $\Sigma$ are transverse,
since the tangent is timelike and both hypersurfaces are spacelike.
For large $T$ its bounded-$v$ deep portion lies below
$(1-\nu_0)T$, and all its lower causal shadows lie in the protected
ball of Equation~\eqref{attach:protected-radius}. Exact earlier entry agreement from Proposition~\ref{prep:signed-entry},
the support conclusion of Proposition~\ref{energy:signed-experiment},
and the attachment therefore identify this entire
neighborhood with the background one. Finite-time smooth geodesic
flow dependence carries the convergent states at $v_0$ through
these collars to seeds $w_T^\pm\to w$. This compact observer
segment is distinct from the packet's moving entry foot near
$v=T$.

The constructed path is a future unit geodesic of the attached
vacuum metric. Equation~\eqref{attach:future-root} also gives
$d\tau/dv=2a/\partial_z\mathcal F_T$ on its late portion, so every
finite experiment has finite proper duration. Its earlier portion
has finite duration by smooth compact flow dependence. At $p_T$
there is an open observation neighborhood in the attached spacetime,
so local geodesic continuation extends it for a positive further
proper time. Proposition~\ref{attach:development} embeds this path
and its continuation into the full MGHD. Its observation time is
therefore strictly below the full-MGHD endpoint. Openness of the
original seed set gives membership of $w_T^\pm$ for large $T$.

\emph{Comparison with the norm in the test.}
The frame in Equation~\eqref{attach:adapted-frame} has the fixed
Lorentz Gram matrix
\[
 g_T^\pm(E_0^\pm,E_1^\pm)=-1,
 \qquad g_T^\pm(E_A^\pm,E_B^\pm)=\delta_{AB},
\]
with all other entries zero. By isometry, the components of
$U_T^\pm$ in it are precisely $(\dot u,p,\Xi^{(A)})$, bounded by
Equation~\eqref{attach:branch-variables}. They belong to a compact
subset of the unit future timelike hyperboloid. It follows that
the positive form
\[
 r_{U_T^\pm}=g_T^\pm
       +2g_T^\pm(\cdot,U_T^\pm)\otimes g_T^\pm(\cdot,U_T^\pm)
\]
is uniformly equivalent at $p_T$ to the Euclidean component form:
it is positive definite for every unit timelike vector, and its
largest and smallest eigenvalues vary continuously on this compact
set. Both this form and the test form are parallel along the
observing geodesic. At $\Sigma$ their comparison is uniform because
$d_T^\pm\to d$ and $w_T^\pm\to w$; equivalently use
Lemma~\ref{test:observer-comparison}. Parallel transport preserves
the comparison constants. This proves the norm assertion and,
by applying it to a vector and its image, the induced operator-norm
assertion. A smooth $e_A$ with the stipulated central direction and
fixed jets of cost $e^{o(T)}$ is obtained by Gram--Schmidt in the
angular chart followed by a constant orthogonal rotation.
\end{proof}
